\section{Layering and axioms: a closed argument must be stratified}
\label{sec:layers_axioms}

To prevent ``interpretation'' from being silently used as ``derivation'', we make the logic auditable by explicitly separating three layers.

\paragraph{Layer 0 (ontological layer).}
Only the language of \emph{states and algebras} is allowed. No external time parameter and no external probability postulate is introduced.

\paragraph{Layer 1 (protocol layer).}
One may choose a scan and readout protocol under finite-resolution constraints. Operational time and statistics are defined \emph{within} the protocol.

\paragraph{Layer 2 (interpretation layer).}
One maps Layer~1 structures to semantic narratives (spacetime, particles, gravity, entanglement). This layer must be marked interpretive and may not be used as a premise of the Layer~0/1 argument.

This stratification is consistent with the HPA--$\Omega$ manuscripts \cite{Ma2025HPA,Ma2025FoP} and is adopted here as a strict audit rule.

Figure~\ref{fig:layers_flow} summarizes the separation and highlights the one-way dependency: Layer~2 interpretation may annotate, but may not be used as a premise.

\begin{figure}[H]
\centering
\includegraphics[width=0.95\linewidth]{fig01_layers_flow}
\caption{Layer separation and the audit rule. Layer~0 uses only states and algebras; Layer~1 specifies scan and finite-resolution readout protocols and induced measures; Layer~2 provides semantic interpretations and is explicitly excluded from the derivation chain.}
\label{fig:layers_flow}
\end{figure}

\subsection{Basic axioms (Omega O1--O4)}
\label{sec:axioms_O1_O4}

\begin{axiom}[O1 (Omega: static global state)]
The universe is specified by a quasi-local operator algebra $\Acal$ together with a unique normalized global state $\omega_\Omega$ on $\Acal$. There is no externally given time-indexed family of states.
\end{axiom}

\begin{axiom}[O2 (finite information)]
For any causally closed region, the effective number of degrees of freedom is bounded by a holographic constraint, typically of the form
\begin{equation}
\dim \hilbert_{\mathrm{region}} \le \exp\!\left(\frac{A}{4\ell_P^2}\right),
\end{equation}
where $A$ is an appropriate boundary area and $\ell_P$ is the Planck length.
\end{axiom}
\begin{remark}
See, e.g., \cite{Bousso2002} for a standard review of holographic bounds and related formulations of the holographic principle.
\end{remark}

\begin{axiom}[O3 (causally local discrete update)]
There exists a discrete-step automorphism $\mathcal{U}:\Acal\to\Acal$ which, in a controlled representation, is implemented by a unitary $U$ and has finite propagation range (a causal locality condition).
\end{axiom}

\begin{axiom}[O4 (holographic map)]
There exists a bulk-to-boundary encoding map $\Phi$ which is approximately isometric on an appropriate code subspace and supports observable transport and reconstruction in the sense of algebraic quantum error correction / entanglement wedge reconstruction.
\end{axiom}
\begin{remark}
For standard constructions and results connecting holography, bulk reconstruction, and quantum error correction, see, e.g., \cite{AlmheiriDongHarlow2015,Pastawski2015,Harlow2017}.
\end{remark}

\subsection{Upgraded axioms: scan--projection readout and Weyl pairs (O5--O6)}
\label{sec:axioms_O5_O6}

\begin{axiom}[O5 (scan--projection readout and induced measures)]
Finite observers obtain ``time'' and ``probability'' from a finite-resolution scan and projection readout of intrinsic phase data. There exists a pointer phase $x\in\RR/\ZZ$ whose scan orbit is sampled at integer ticks $t\in\ZZ_{\ge 0}$ by
\begin{equation}
x_t = x_0 + t\alpha \pmod 1,\qquad \alpha\notin\QQ.
\end{equation}
Finite-resolution readout is described by a family of effects $\{E_k^{(\varepsilon)}\}$ (resolution parameter $\varepsilon$) inducing probabilities
\begin{equation}
P_k^{(\varepsilon)} = \omega_{\mathrm{eff}}\!\left(E_k^{(\varepsilon)}\right),\qquad \sum_k E_k^{(\varepsilon)}=\ind,
\end{equation}
where $\omega_{\mathrm{eff}}$ denotes the effective state in the observer sector.
\end{axiom}

\begin{remark}[standard measurement-theory interface]
The family $\{E_k^{(\varepsilon)}\}$ is a POVM in the usual sense (effects summing to the identity), and $E\mapsto \omega_{\mathrm{eff}}(E)$ is a state on the corresponding effect algebra; see \cite{Kraus1983,NielsenChuang2000} for standard measurement-theory foundations. In finite-dimensional Hilbert-space settings, Gleason-type theorems constrain such probability assignments and justify the Born form under standard assumptions; see \cite{Gleason1957,Busch2003} for classical statements and extensions.
\par
In algebraic formulations, one may take $\omega_{\mathrm{eff}}$ as the restriction of the global state $\omega_\Omega$ to an observer-accessible subalgebra (or via a conditional expectation onto such a subalgebra), which recovers the familiar reduced-state/partial-trace picture when a tensor-factorization is available; see \cite{BratteliRobinson1997,Haag1996}.
\end{remark}

\begin{remark}[continuum parameters versus finite information]
Although O5 models $x\in\RR/\ZZ$ and $\alpha\notin\QQ$ as continuum-valued parameters, the framework does not assume that an observer can access their infinite precision. Operationally, access is only through finite-resolution effects $\{E_k^{(\varepsilon)}\}$, and O2 imposes a finite effective information bound. Continuum phase variables should therefore be read as ideal coordinates for a protocol limit, not as physically available infinite storage.
\end{remark}

\begin{axiom}[O6 (unitary scan algebra: a Weyl pair)]
On an effective observer Hilbert space $\hilbert_{\mathrm{eff}}$, there exist a scan unitary $U_{\mathrm{scan}}$ and a conjugate phase unitary $V$ satisfying the Weyl relation
\begin{equation}
U_{\mathrm{scan}} V = \e^{2\pi \iu \alpha} V U_{\mathrm{scan}}.
\end{equation}
A canonical covariant model is given on $L^2(\RR/\ZZ)$ by
\begin{equation}
(U_{\mathrm{scan}}\psi)(x)=\psi(x+\alpha),\qquad (V\psi)(x)=\e^{2\pi \iu x}\psi(x).
\end{equation}
\end{axiom}

\begin{assumption}[R1 (orbit regularization / finite part)]
Regulated-to-continuum passages along scan orbits are fixed by a canonical regularization convention: ``Abel first, then limit'', which selects a unique finite part for orbit traces and related divergent sums.
\end{assumption}
\begin{remark}
The Abel-regularization convention is a standard summability prescription: one introduces an Abel damping parameter (e.g.\ $r^t$ with $0<r<1$), performs the relevant summation at fixed $r$, and then takes $r\uparrow 1$ to select a canonical finite part when it exists. See \cite{Hardy1949} for classical background on Abel summability and related regularization methods.
\end{remark}
\begin{remark}
Appendix~\ref{app:abel_regularization} records an explicit Abel-mean lemma for the irrational-rotation orbit of O5, together with a concrete ``finite part'' formula on trigonometric polynomials.
\end{remark}

The purpose of O5--O6 is to push ``probability'' and ``measurement'' down from Layer~2 to Layer~1: probabilities are induced by finite-resolution effects rather than postulated as external sampling semantics.


